\chapter{Singular limits and deep horns}
\label{ch:singular-limits-v4}

At a singular time, bounded scalar curvature singles out a region that
survives smoothly. Its complement need not be a finite collection of points,
and the surviving metric need not be complete. We construct this region in
the original spacetime and then understand its ends. Far enough along an end,
a neck of fixed accuracy improves to a strong neck of any prescribed
accuracy. The improvement includes its backward flow, which is what a
subsequent surgery construction needs.

Throughout this chapter the dimension is three, $R$ denotes scalar curvature,
and $\operatorname{Rm}$ denotes the full curvature tensor. A strong neck
includes an actual backward flow cylinder and normalized metric estimates,
not just a cylindrical terminal metric. We use the neck and cap definitions
and their four canonical alternatives from
Section~\ref{sec:neck-cap-definitions}.
The analytic inputs are curvature evolution and local derivative estimates
(Proposition~\ref{prop:rf-evolution} and Theorem~\ref{thm:rf-shi}),
the pointed compactness construction of Section~\ref{sec:rf-compactness},
the differential Harnack inequality (Theorem~\ref{thm:rf-harnack}), and
the stated noncollapsing hypotheses on the original flows.
The two-dimensional ancient classification enters only at the final horn
improvement. A finite backward limit is not yet an ancient round cylinder.

\section{Bounded distance and the cone obstruction}
\label{sec:singular-bounded-distance}

Write $\nu\geq0$ for the negative part of the least sectional curvature.
The weak pinching condition used here is, on every included slice,
\begin{equation}
 R\geq-6,\qquad
 \nu>0\ \Longrightarrow\ R\geq2\nu(\log\nu-3).
 \label{eq:singular-weak-pinching}
\end{equation}
The normalized Hamilton--Ivey inequalities imply this condition on
nonnegative times. Nonnegative curvature also implies it, without imposing
a restriction on the time origin.

\begin{theorem}[Bounded curvature at bounded distance]
\label{thm:singular-bounded-distance}
\uses{thm:neck-cap-cover}
\lean{PoincareMT.m28BoundedDistance}
\uses{lem:singular-cone-obstruction,thm:rf-shi,thm:rf-harnack}
There exists $0<\varepsilon_0\leq1/200$ with the following property.
For every $0<\varepsilon\leq\varepsilon_0$, $C>0$, and $A\geq0$ there
are $D_0,D>0$ such that the following holds for every generalized flow $F$
satisfying \eqref{eq:singular-weak-pinching}. Let $t$ be an included time,
$x\in F_t$, and $Q=R(x,t)\geq D_0$. Suppose every point of the time-$t$
slice with scalar curvature at least $4Q$ has a strong
$(\varepsilon,C)$ canonical neighborhood. Then
\begin{equation}
 y\in B_{g(t)}(x,AQ^{-1/2})\quad\Longrightarrow\quad R(y,t)\leq DQ.
 \label{eq:bounded-distance}
\end{equation}
There is also a choice of $D_0,D$ for the following replacement of the
same-time hypothesis: for every included $s\leq t$ and every $a<s$ there
is an included $u\in(a,s]$ whose entire slice has strong canonical
control above the fixed threshold $4Q$.
\source{MT2007}\dcref{Theorem 10.2, pp. 245--265}
\end{theorem}

The constants in the two assertions need not be identical. The accuracy
threshold precedes $\varepsilon,C,A$; the curvature constants precede the
flow and the selected points. In the dense-time assertion the good time
$u$ precedes every spatial point on that slice, and may equal $s$.

\subsection{The two successive limits}

Suppose the same-time estimate fails. Choosing counterexamples with both
proposed constants tending to infinity gives points $x_n,y_n$ with bounded
normalized distance and
\[
 Q_n=R(x_n,t_n)\longrightarrow\infty,\qquad
 R(y_n,t_n)/Q_n\longrightarrow\infty.
\]
On a path from $x_n$ to $y_n$, take its last
crossing of an intermediate scalar level. Its remaining segment lies in the
high-curvature region, where the canonical alternatives apply.

Neck overlaps order the central spheres along this segment. If a cap occurs,
its core is retained and the outgoing end is joined to the neck chain.
The alternatives are a tube, a capped tube, or a compact component. Two cap
cores covering the segment would bound the ratio of the endpoint scalars,
so cannot persist in this sequence. A minimizing sequence of paths between
the selected endpoints stays in a compact part of the tube: disjoint central
spheres enclose a compact cylinder; when a cap is present its compact core
replaces one end. The neck exit estimates prevent an excursion into the
remaining end from shortening the path. Thus one obtains intrinsic
minimizers without assuming completeness of the high-curvature region
\cite[Claims 10.3--10.6, pp. 248--251]{MT2007}.

Choose the critical radius at which uniform curvature control breaks down.
The choice is made after fixing the source sequence and the first retained
neck: it is the supremum of radii for which this family of source necks has
the required curvature, distance, and backward-cylinder bounds. On strictly
smaller balls the canonical estimates give curvature and derivative bounds,
and neck charts give the local volume needed for compactness. The closed
three-quarter part of the first neck is captured in one fixed exhaustion
stage for every sufficiently late source index and every point of that core.
Exhaust those balls and extract compatible
pointed limits. Pinching makes the limiting sectional curvature nonnegative:
under a diverging scalar normalization, a fixed positive negative-curvature
defect would contradict the logarithm in
\eqref{eq:singular-weak-pinching}.

At the critical boundary there are points of arbitrarily large scalar. The
retained source family gives a positive-side neck cover. Its total volume is
finite: the source critical regions have a common volume bound, and metric
convergence transfers that bound to every compact exhaustion stage, then to
their union. A bound $R\leq B$ at neck centers gives a common scale floor.
The balls of radius $\delta=B^{-1/2}/16$ about those centers have a fixed
positive volume lower bound, and their doubled balls have compact closure
inside the neck charts. A $2\delta$-separated set of such centers has
bounded cardinality by disjoint-ball volume comparison. A maximal such set
therefore gives a finite cover by doubled balls. Thus the bounded-scalar
centers lie in one compact set, without assuming ambient completeness.

The selected positive side has noncompact closure, so its scalar is
unbounded. The neck-chain construction selects a separating section within
that side. Its outgoing half has closure in the ambient tube, where a
continuous height takes values below one. For each scalar bound, compactness
of the bounded-scalar centers gives a maximum height strictly below one.
Every later tail therefore exceeds that scalar bound. Recut the selected
section inside this same ambient tube. In
coordinates $s\in(-\varepsilon^{-1},\varepsilon^{-1})$ one may take the
recut collar between $s=-\varepsilon^{-1}/4$ and
$s=\varepsilon^{-1}/4$; the initial separating sphere, source paths, and
limit metric are unchanged. The recut has a connected open tube $U$ with
compact boundary, finite intrinsic diameter, positive scalar curvature, and
an end along which $R\to\infty$. Its proper height has uniform scalar
divergence, so every sufficiently far tail is in the same positive side and
the limiting strong-neck cylinders provide backward flow near its points.
This first limit is incomplete. No completeness of the limit has been used.

In the intrinsic metric completion of $U$, the escaping end is a missing
point $E$. The neck diameter and separation estimates imply that a
sufficiently deep tail has arbitrarily small diameter; hence sequences
escaping that tail define the same completion point. Parameterize radial
segments by their distance $r$ from $E$. Nonnegative sectional curvature,
applied to confined minimizing segments, gives limits of the normalized
chords between radial segments. If $\xi,\eta$ denote their germs, the
resulting cone distance has the form
\begin{equation}
 d_{\rm cone}((\xi,r),(\eta,s))^2
       =(r-s)^2+rs\,d_L(\xi,\eta)^2.
 \label{eq:singular-cone-chord}
\end{equation}
Here $d_L$ is a chord metric on the space of germs, obtained by quotienting
zero limiting chord distance. For each pair of rays, finite-segment
comparison gives a chord-defect limit $K(P,Q)\in[0,4]$, and the bound holds
simultaneously on every strict subrectangle of their lengths. The limit is
joint as both radii tend to zero. Finite nets on a compact annulus then make
the convergence uniform there; passing to a further finite net is the only
diagonal step. The neck spheres give at least two distinct germs. The
construction retains the intrinsic metric of this same tube; it does not
replace an incomplete limit by its completion and invoke uniqueness of
complete pointed limits. In the eventual source chart the radial buffer is
quantitative: a closed ball of radius $a/64$ is mapped into the open ball of
radius $a$, and the smaller ball of radius $a/128$ is reserved for the
comparison. This leaves a genuine open annulus away from the cone vertex on
which the source distances and the curvature-scale distances are the same
data
\cite[Sections 10.3--10.4, pp. 253--263]{MT2007}.

There are two rescalings near $E$: distance to the missing point and scalar
curvature. Their comparability is indispensable. On one fixed radial
segment choose $z_i$ with $d_i=d(z_i,E)\downarrow0$. The centered neck
around $z_i$ cannot meet the missing point, giving a positive lower bound
for $R(z_i)d_i^2$. If this product diverged, its central neck sphere after
distance normalization would have diameter tending to zero while separating
two fixed radial annuli. Its limit would make a point at positive radius
separate those annuli of the cone in \eqref{eq:singular-cone-chord}.
Two distinct ray germs exclude that possibility. The construction yields
a subsequence with
\begin{equation}
 \frac1{16}<R(z_i)d_i^2\leq B,
 \qquad R(z_i)d_i^2\longrightarrow\lambda>0.
 \label{eq:singular-scalar-radius}
\end{equation}
The numerical lower bound is weaker than the printed neck estimate; only
strict positivity is needed.

Rescale by $R(z_i)$ and retain the actual strong-neck backward cylinders.
There is a smooth local backward limit with base scalar one.
Equation~\eqref{eq:singular-scalar-radius} keeps the basepoint a positive
finite distance from the cone point. The same-source distance convergence
identifies an open neighborhood of the terminal slice isometrically with
an annulus of \eqref{eq:singular-cone-chord}. The comparison uses the same
points and a further strict subsequence of the original sequence. It is a
local identification away from the vertex
\cite[Claim 10.31 and Sections 10.5--10.6, pp. 263--265]{MT2007}.

\subsection{Why the local cone is impossible}

\begin{lemma}[Finite backward cone obstruction]
\label{lem:singular-cone-obstruction}
\lean{PoincareMT.M28.no_positive_terminal_scalar_of_chord_annulus_embedding}
\uses{prop:rf-evolution}
Let $g(t)$, $a\leq t\leq b$, $a<b$, be a smooth three-dimensional Ricci
flow with nonnegative curvature operator. If an open part of its terminal
metric is isometric to an open part of a metric cone away from the vertex,
then its terminal scalar curvature vanishes there. Neither completeness nor
an ancient extension is required.
\end{lemma}

\begin{proof}
Pull back the radial potential $f=r^2/2$ to the terminal open region.
The cone chord identity gives, along every short affinely parameterized
terminal geodesic $\gamma$,
\[
 f(\gamma(s))=(1-s)f(\gamma(0))+sf(\gamma(1))
                  -\tfrac12s(1-s)|\dot\gamma(0)|^2.
\]
Initially $f$ is locally Lipschitz, since radius is Lipschitz and bounded
on a smaller annulus. Joint smooth dependence of geodesics and the inverse
function theorem give normal neighborhoods uniform in nearby centers.
Rademacher's theorem supplies a differentiability center $q$ whose normal
neighborhood contains any prescribed point. The quadratic identity gives
\[
 f(\exp_qv)=f(q)+df_q(v)+\tfrac12|v|^2
\]
there. Its right side is smooth, proving smoothness near the prescribed
point. Differentiating along geodesics now gives $\nabla^2f=g(b)$.
Thus $Z=\nabla f$ satisfies
$\nabla_vZ=v$. Commuting derivatives gives
$\operatorname{Rm}(v,w)Z=0$, hence $\operatorname{Ric}(Z,Z)=0$.

Use the terminal null-section consequence of curvature evolution: on a
nonnegatively curved flow with a nontrivial included backward interval, a
smooth field Ricci-null on an open terminal region satisfies
\[
 \operatorname{Ric}_b(\nabla_vZ,\nabla_wZ)=0.
\]
Here is the calculation. Positive semidefiniteness makes
$\operatorname{Ric}_b(Z,W)=0$ for every $W$. Differentiate this identity
twice, with $e_i$ an orthonormal basis, to obtain the tensor Laplacian
\[
 (\Delta\operatorname{Ric}_b)(Z,Z)
       =2\sum_i\operatorname{Ric}_b(\nabla_{e_i}Z,\nabla_{e_i}Z).
\]
The Ricci reaction term vanishes on the null direction under nonnegative
sectional curvature. At a fixed point the function
$t\mapsto\operatorname{Ric}_t(Z,Z)$ is nonnegative and zero at $b$,
so its left derivative is nonpositive. The Ricci evolution equation
forces the displayed sum of nonnegative terms to be zero.
Positive semidefiniteness and linearity extend this nullity from the basis
to every $\nabla_vZ$. This proves the stated null-section consequence
and the local maximum-principle obstruction of
\cite[Proposition 4.22, p. 74]{MT2007}.
Since $\nabla Z$ is the identity, terminal Ricci and scalar curvature vanish.
\end{proof}

The smooth limit above has scalar one at its basepoint, a contradiction.
This proves the same-time assertion. Strong canonical neighborhoods were
needed because the obstruction concerns a terminal cone in a backward flow.

\subsection{Inclusive dense times and path rebasing}

Let $D_0,D$ be the same-time constants for $A'=(A+1)\sqrt2$. Use
$D'_0=D_0/2$ and $D'=\max\{4,2D\}$. Suppose the dense-time conclusion
fails at $(x,t)$, with scalar $Q\geq D'_0$, at a point $y$.
A piecewise smooth path from $x$ to $y$ has length less than
$(A+1)Q^{-1/2}$. Its compact image is covered by finitely many flow boxes.
On a sufficiently close included time their compatible transports preserve
the length bound and the strict inequalities
$R(x_u,u)<2Q$, $R(y_u,u)>D'Q$. Choose one good time $u$ for this whole path.

Scalar continuity on the path gives a last crossing $z$ of $2Q$. Its tail
stays above $2Q$ and has length less than
\[
 (A+1)Q^{-1/2}=A'(2Q)^{-1/2}.
\]
The good slice supplies canonical neighborhoods above $4Q$, hence above
$4R(z,u)=8Q$, and $R(z,u)=2Q\geq D_0$. The same-time theorem based at
$z$ gives $R(y_u,u)\leq2DQ$, a contradiction. At an included left
endpoint the good time may be that endpoint itself. No spatially varying
choice of good times is used.

\section{Controlled generalized-flow limits}
\label{sec:singular-blowup}

\begin{definition}[Scalar-normalized sequence and common controls]
\label{def:singular-blowup-sequence}
\lean{PoincareMT.GeneralizedBlowupSequence,PoincareMT.M30CommonBlowupControls}
A blow-up sequence consists of generalized flows $F_k$, included times
$t_k$, points $x_k\in(F_k)_{t_k}$, and $Q_k=R(x_k,t_k)>0$ tending to
infinity. Its normalized time is $s=Q_k(t-t_k)$ and metric $Q_kg(t)$.
Write $B_k(A)=B_{g(t_k)}(x_k,A/\sqrt{Q_k})$.

The common controls are positive constants
$\varepsilon,C,\kappa,r_0,\mu,C_a$ with these properties. Each flow
is either Hamilton--Ivey pinched on nonnegative times or nonnegatively
curved. The inclusive dense-time condition holds at every earlier included
time with cutoff $4Q_k$. Above that cutoff at earlier times,
\[
 |\nabla R|\leq C_aR^{3/2},\qquad |\partial_tR|\leq C_aR^2,
\]
where the temporal derivative is in actual flow boxes. For each $A>0$,
eventually $B_k(A)$ has compact closure and every point in it is
$\kappa$-noncollapsed on physical scales at most $r_0$. The maximal
backward worldline of every such point $y$ includes the normalized interval
\begin{equation}
 \left[-\frac{\mu Q_k}{\max\{Q_k,R(y,t_k)\}},0\right].
 \label{eq:singular-survival}
\end{equation}
Maximality is among compatible worldlines in the actual spacetime.
\end{definition}

\begin{proposition}[Bounded distance on a blow-up sequence]
\label{prop:singular-m29}
\lean{PoincareMT.m29GeneralizedBoundedDistance}
\uses{def:singular-blowup-sequence,thm:singular-bounded-distance}
Below the universal accuracy threshold, for each fixed $A>0$ there is
$D>0$, depending only on $\varepsilon,C,A$, such that eventually
$R(y,t_k)\leq DQ_k$ for every $y\in B_k(A)$.
\end{proposition}

\begin{proof}
Use the dense-time constants of
Theorem~\ref{thm:singular-bounded-distance}. Both pinching branches imply
\eqref{eq:singular-weak-pinching}; eventually $Q_k\geq D_0$.
Since $d_{Q_kg}=\sqrt{Q_k}\,d_g$, the normalized radius-$A$ ball is
exactly $B_k(A)$. Apply the estimate there. The gradient, survival and
noncollapse controls are unnecessary for this proposition.
\end{proof}

For a positive finite or infinite horizon $T_0$, define
\begin{equation}
 J_{T_0}=\{s\leq0:-s<T_0\}
 =\begin{cases}(-T_0,0],&T_0<\infty,\\(-\infty,0],&T_0=\infty.
 \end{cases}
 \label{eq:singular-limit-domain}
\end{equation}
Long controls additionally require that, for every $0<T<T_0$ and $A>0$,
eventually there is an actual compatible cylinder on
$(-T,0]\times B_k(A)$, equal to the identity at zero, at every image
point of which the flow is $\kappa$-noncollapsed on physical scales at
most $r_0$. The threshold in $k$ may depend on $T,A$.

\begin{definition}[The convergence retained in the limit]
\label{def:singular-convergence}
\lean{PoincareMT.GeneralizedBlowupConvergence,PoincareMT.BlowupLimitFlow}
\uses{def:singular-blowup-sequence}
Convergence on a backward interval $J$ consists of one strict subsequence,
a connected pointed smooth three-manifold $(N,p)$, a Ricci flow $g_\infty$
on $J$, and increasing connected relatively compact open exhaustions
$U_j\subset N$. Positive times $\tau_j$ have
$[-\tau_j,0]\subset J$ and are cofinal on compact subsets of $J$.
Compatible embeddings
\[
 [-\tau_j,0]\times U_j\longrightarrow F_{k_j},\qquad
       (s,z)\longmapsto\Psi_j(s,z)
\]
have physical time $t_{k_j}+s/Q_{k_j}$ and send $(0,p)$ to $x_{k_j}$.
Their pulled-back normalized metrics converge with every mixed
space--time derivative on compact subcylinders, including time zero.
Boundary derivatives are within the stated domains. For every $A>0$,
the zero-time image eventually contains the entire original ball
$B_{k_j}(A)$.

Every included limit slice is complete, has nonnegative curvature operator,
and $R_\infty(p,0)=1$. For each compact $I\subset J$ there is $B_I<\infty$
such that $|\operatorname{Rm}_\infty|\leq B_I$ on all of $N\times I$.
\end{definition}

\begin{theorem}[Short and long controlled limits]
\label{thm:singular-m30}
\lean{PoincareMT.m30ControlledGeneralizedBlowupLimits}
\uses{prop:singular-m29,def:singular-convergence,thm:rf-shi,thm:rf-harnack}
One can choose $0<\varepsilon_1\leq1/400$ before the sequence and its
constants such that common controls with $\varepsilon\leq\varepsilon_1$
give convergence as in Definition~\ref{def:singular-convergence} on some
closed $[-\tau,0]$, $\tau>0$. Long controls give such convergence on
$J_{T_0}$. The long limit is $\kappa$-noncollapsed on scales at most
$r_0$. If $T_0=\infty$, it is an ancient $\kappa$-solution on this
same carrier, with these same metrics and Levi--Civita connections.
\source{MT2007}\dcref{Theorems 11.1 and 11.8, pp. 267--279}
\end{theorem}

The finite-scale assertion tests $r\leq r_0$ and included intervals
$(s-r^2,s]\subset J_{T_0}$. If full curvature is at most $r^{-2}$ on
the fixed terminal ball $B_{g_\infty(s)}(z,r)$ at all these times, its
time-$s$ volume is at least $\kappa r^3$.

\subsection{A common short lifetime}

Proposition~\ref{prop:singular-m29} initially gives a bound depending on
$A$. Above the scalar guard, integrate
$|\partial_t(R^{-1})|\leq C_a$ to obtain a backward interval where
scalar stays within a fixed multiple of its terminal value. Below the
guard, start the estimate at a first crossing of the guard. The spatial
bound controls variation on curvature-scale balls. With
\eqref{eq:singular-survival}, these estimates supply bounded cylinders
over each fixed $B_k(A)$. Local derivative estimates and the volume bound
give pointed compactness first for terminal metrics and then for local
backward flows on an exhaustion.

All-radius source coverage and compactness of original bounded balls
ensure completeness of the terminal metric. More precisely, for a limit ball
of radius $A$ choose a source ball of radius $A+1$ and then a compact
exhaustion stage which contains its closure. The convergence maps cover the
smaller limit ball and preserve a strict distance margin to the boundary of
the source ball. A finite-length curve of length at most $A$ therefore has
its initial segment in that stage and cannot leave the source ball while its
continuation is transferred. The compact closure supplies a convergent
endpoint. Increasing $A$ gives every bounded limit ball, so a finite-length
Cauchy sequence converges. This is a global endpoint argument, not a local
coordinate assertion; completeness is required at zero, not only in the
interior.

To obtain a common lifetime, first transfer canonical neighborhoods at
terminal points of scalar greater than four. Capture each selected canonical carrier in
a buffered compact set, sample original certificates at good times, and
pass the normalized jets to the limit. A single enlargement gives accuracy
$2\varepsilon$ and constant $2C$. In particular a cap supplies a neck
whose center scalar is comparable with that of the point. Every
high-curvature point is thus related to such a neck, unless the component
is compact.

The terminal scalar is globally bounded. For a compact component this is
continuity. On a complete noncompact strictly positively curved component,
the neck-scale lower bound excludes arbitrarily small necks, while the
preceding scalar comparison would produce them from unbounded scalar.
Here we use the small-neck scale argument in the proof following
Lemma~\ref{lem:ancient-terminal-cylinder}: its point-soul, nested-side
and flux construction applies to every complete strictly positively curved
three-manifold, independently of a soliton equation.
If a curvature plane is null, the local backward flows and tensor maximum
principle yield parallel null directions. These local splittings agree on
connected overlap regions containing a fixed nonflat basepoint. Together
with one related neck they bound curvature in the remaining directions.
This is the dichotomy in \cite[Claim 11.7, pp. 270--271]{MT2007}.
The local backward flows used at this stage need not have a common lifetime.

Let $B\geq4$ be the terminal scalar bound. Each fixed source ball eventually
has normalized terminal scalar at most $B+1$. The scalar ODE and
maximal-worldline condition now give a common positive backward lifetime,
depending on $B,C_a,\mu$ before the radius is chosen. Extract there and
restrict to a slightly shorter closed interval $[-\tau,0]$. Uniform
metric comparison transfers terminal completeness to its other times.
A diagonal choice of space, time and derivative order gives one set of
embeddings for all mixed jets.

\subsection{Finite left endpoints and the whole horizon}

An extracted interval might appear to stop at a finite $-T$ strictly
inside the prescribed horizon. We first construct a smooth complete metric
at that endpoint, before invoking the endpoint curvature estimate.
Let $B$ bound scalar curvature on the already included time-zero slice.
Differential Harnack on the interval $(-T,0]$ gives
\[
 R(z,s)\leq\frac{BT}{T+s}\qquad(-T<s\leq0).
\]
The distance variation inequality integrates the square root of this
bound. Although curvature may appear to diverge as $s\downarrow-T$,
its square root is integrable. Distances between fixed points therefore
increase backward by a bounded amount depending on $B,T$. If a compact
set contains at each time a point of uniformly bounded scalar, bounded
distance propagates that bound across the set. Apply the estimate on
approximating source slices, where the original canonical thresholds are
available, then pass to the limit
\cite[Lemma 11.11 and Claim 11.12, pp. 273--274]{MT2007}.

On a compact limit scalar minimum monotonicity supplies this anchor.
On a noncompact limit choose a distant $y$ and, by comparison along rays,
a point $z$ with $d_0(y,z)=d_0(p,y)$ and
$d_0(p,z)>3d_0(p,y)/2$. If $R(y,s)$ became unbounded toward $-T$,
its canonical neighborhood would have vanishing scale. It cannot be a
cap, and minimizing paths toward $p$ and $z$ cannot exit the same neck
end: small exit angle and Toponogov comparison would contradict terminal
separation even allowing the bounded backward distance change. They exit
opposite ends. The central sphere then separates $p,z$ inside a ball of
radius tending to zero at $y$ in the terminal metric. This is impossible:
a fixed path avoiding $y$ stays a positive distance from it. Distant
points consequently remain bounded and anchor every fixed compact set
\cite[Lemma 11.13--Corollary 11.16, pp. 275--277]{MT2007}.

The resulting compact-local bounds up to $-T$ give a smooth left-endpoint
metric on every captured compact set. Nonnegative Ricci curvature makes this
backward metric dominate the complete zero-time metric. An endpoint-metric
Cauchy sequence therefore converges in the zero metric and then, by local
equivalence on a captured ball, in the endpoint metric. The endpoint is
complete. The endpoint curvature theorem then upgrades the compact-local
bounds to one global scalar bound; the guarded ODE and the surviving source
cylinders extend convergence by one common positive duration. Harnack
supplies a spatially global bound on every closed slab of the extended
interval
\cite[Claim 11.17, pp. 277--278]{MT2007}.

Take the supremum of attainable horizons, retaining the original sequence
and source embeddings. The first short construction gives one terminal
scalar bound $Q$ for every later finite limit. Choose increasing horizons
$T_j$ cofinal below the supremum. Harnack bounds scalar on $[-T_j,0]$ in
any limit of duration at least $T_{j+1}$ by
\[
 Q\,\frac{T_{j+1}}{T_{j+1}-T_j}.
\]
Thus each strict prefix has a coefficient fixed before choosing a longer
limit or a spatial radius. For any finite list of prefixes, take one limit
longer than the list and transfer all these bounds to its retained original
sources. Diagonal extraction from these source cylinders, with the uniform
terminal volume bound, gives one limit and compatible embeddings on every
compact cylinder. It does not require unrelated finite limits to have been
identified in advance. If the supremum were
strictly below $T_0$, the complete endpoint and global bound just obtained
would extend it by a common duration, a contradiction. Thus convergence is
on $J_{T_0}$, with no added slice at a finite excluded $-T_0$.

Noncollapse passes through the same metric and volume convergence.
Allow a strict curvature and radius margin on a fixed tested parabolic
ball, apply the source inequality, then let the margin decrease to zero.
At infinite horizon every finite normalized scale is eventually below
$\sqrt{Q_k}\,r_0$ in the source. Thus the ancient limit is noncollapsed
at all scales. Nonnegative curvature, completeness and $R(p,0)=1$ give
its remaining defining properties; Harnack bounds all past scalar by the
terminal global bound. This identifies an ancient solution without yet
classifying it.

There is also a geometric branch used in horn continuation. It assumes
compact normalized terminal balls, a terminal volume bound
$\operatorname{Vol}(B_k(\rho_*))\geq v_*Q_k^{-3/2}$ for fixed
$\rho_*,v_*>0$, and, for each $T<T_0$, a constant $B_T$ chosen before
$A,\eta>0$ such that eventual actual cylinders on $[-T,0]\times B_k(A)$
have $|\operatorname{Rm}|\leq B_TQ_k$ and negative curvature part at
most $\eta Q_k$. Extraction gives
Definition~\ref{def:singular-convergence} on $J_{T_0}$. This branch asserts
neither noncollapse nor an ancient-solution certificate.

\section{The retained region at a singular time}
\label{sec:singular-retained}

Fix finite $T$ and an old interval exactly $[0,T)$. The discrete set of
singular times includes $T$; regular old slices are compact. Choose
$t_-<T$ in the old interval and a reference carrier $M$ with compatible
diffeomorphisms $\Phi_t:M\to F_t$ for $t\in[t_-,T)$. Their metric
pullbacks form a Ricci flow $g(t)$, and their spacetime map identifies
exactly this time window and agrees locally with the vertical flow.
In particular $M$ is compact.

\begin{definition}[Singular-time data and retained region]
\label{def:singular-retained-region}
\lean{PoincareMT.SingularTimeAssumptions,PoincareMT.SingularTimeReference.regularLimitSet}
Assume the normalized pinching inequalities
\[
 R\geq-\frac6{1+4t},\qquad
 \nu>0\ \Longrightarrow\ R\geq2\nu(\log\nu+\log(1+t)-3),
\]
and a global scalar lower bound. Fix $r_0,C,C_a>0$, $0<\varepsilon<1/4$.
At all old points with $R\geq r_0^{-2}$ assume
\begin{equation}
 |\partial_tR|\leq C_aR^2,\qquad |\nabla R|\leq C_aR^{3/2}.
 \label{eq:singular-analytic-bounds}
\end{equation}
At zero and every nonsingular old time assume strong
$(\varepsilon,C)$ canonical neighborhoods above that threshold.
Put $c_T=10^{14}$ and require $\varepsilon_T=c_T\varepsilon\leq1/200$.
With $R_{\rm ref}(t,x)=R(\Phi_t(x),t)$, define
\begin{equation}
 \begin{split}
 \Omega=\{x\in M:\ &\exists B\in\mathbb R\ \forall t_0<T\
 \exists t\in[t_-,T),\\
 &t_0<t\text{ and }R_{\rm ref}(t,x)\leq B\}.
 \end{split}
 \label{eq:singular-retained-region}
\end{equation}
\end{definition}

The definition uses a finite liminf. Eventual boundedness is a conclusion.
At exceptional old times we use the all-time analytic inequalities and
continuity, rather than assume canonical control.

\begin{definition}[Terminal extension and strong horn]
\label{def:singular-terminal-horn}
\lean{PoincareMT.GeneralizedFlowExtension,PoincareMT.StrongHorn}
\uses{def:singular-retained-region}
A terminal extension retains the old flow through an open spacetime
embedding, with exact old metric and scalar pullbacks, and adds only time
$T$ if its terminal carrier is nonempty. A strong $\alpha$-horn is a
proper smooth parameterization $S^2\times[0,1)\to E_T$, with a smooth
collar across its boundary sphere, whose image is a carrier $H$.
Every point of $H$ centers a terminal strong $\alpha$-neck in the extended
flow; the boundary sphere is itself a central sphere of one such neck.
These necks need not be contained in $H$.
\end{definition}

\begin{theorem}[Regular terminal limit]
\label{thm:singular-m31}
\uses{thm:neck-cap-cover}
\lean{PoincareMT.m31SingularRegularLimit}
\uses{def:singular-retained-region,def:singular-terminal-horn,prop:rf-evolution,thm:rf-shi}
Given the neck--cap topology theorem with threshold
$\varepsilon_{\rm top}>0$, there is $\varepsilon_0>0$ with
$c_T\varepsilon_0\leq\varepsilon_{\rm top}$ such that every
singular-time datum with $\varepsilon\leq\varepsilon_0$ has the
following conclusion. The set $\Omega$ is open. A terminal extension
$E$ has a smooth open embedding $\iota:E_T\to M$ with image exactly
$\Omega$. The metrics $\iota^*g(t)$ converge in every spatial derivative
on compact subsets to the actual terminal metric. Its scalar is the
Levi--Civita scalar, is bounded below, and is proper. Time $T$ is included
exactly when $\Omega$ is nonempty. The compatible terminal gluing map
has domain $(t_-,T]\times E_T$; its image and the old spacetime cover $E$.

Every terminal point with $R>r_0^{-2}$ has strong canonical control with
accuracy $\varepsilon_T=c_T\varepsilon$ and constant $2C$.
Each end of each terminal component has a tail contained in a strong
$\varepsilon_T$-horn.
\source{MT2007}\dcref{Theorem 11.19 and Lemmas 11.21--11.30, pp. 279--286}
\end{theorem}

An end is represented by nested noncompact components of the complements
of a compact exhaustion. No finiteness of terminal components or horns
is asserted.

\section{Terminal metric and neighborhoods}
\label{sec:singular-terminal}

\subsection{From a liminf to smooth convergence}

At a retained point choose a late time with $R\leq B$. Above the guard,
\eqref{eq:singular-analytic-bounds} gives
$|\partial_t(R^{-1})|\leq C_a$. Choose that time so close to $T$ that
scalar beginning below a fixed number larger than
$\max\{B,r_0^{-2}\}$ cannot double before $T$. If it first crosses
the guard, start the estimate there. This proves an eventual bound along
the whole remaining worldline. On the positive high-curvature region,
\[
                  |\nabla R^{-1/2}|\leq C_a/2.
\]
A small ball at that late time therefore has a uniform scalar bound;
the temporal argument gives a bound on its remaining cylinder. The ball
lies in $\Omega$, proving openness.

Pinching converts scalar bounds into full curvature bounds. Interior
local derivative estimates on a smaller cylinder bound every curvature
derivative up to $T$. A compact subset of $\Omega$ needs only finitely
many such cylinders. In fixed coordinates the Ricci equation bounds every
spatial derivative of $\partial_tg$. Integration from $s$ to $t$ makes
each metric jet Cauchy as $s,t\uparrow T$. Thus the entire family
converges. Two-sided metric comparison keeps the limit positive definite,
and compatibility of coordinate limits gives a smooth metric $g(T)$.

The formulas for its Levi--Civita connection and curvature are continuous
in a positive definite metric and its derivatives. Its scalar is therefore
the limit of reference scalars. Passing to the limit in the time integral
of $\partial_tg=-2\operatorname{Ric}$ proves the Ricci equation with
left derivative at $T$. This is an actual terminal flow slice.

For properness, take a sequence in $\Omega$ with bounded terminal scalar.
Compactness of $M$ gives a subsequential limit $x\in M$. The temporal
estimate backward gives one time interval and one scalar bound for the
sequence; the gradient estimate gives one spatial radius at its initial
time. Eventually $x$ lies in those bounded cylinders, hence is in
$\Omega$. Scalar continuity now makes every closed bounded scalar
sublevel compact in $\Omega$. The scalar lower bound completes properness
as a map to $\mathbb R$ \cite[Lemma 11.21, pp. 280--281]{MT2007}.

Use $\Omega$ with this metric as terminal carrier, transported to the
chosen smooth model $E_T$. Glue its half-closed cylinder to the old
spacetime by the reference maps. Their vertical compatibility makes the
identifications agree in flow boxes; metrics and time vector fields agree
on the overlap. The old embedding is open, the terminal gluing map is open
in the relative time topology, and only terminal points are added.
If $\Omega$ is empty, retain the old flow and use an empty terminal
carrier. All identities remain valid, with $T$ excluded from the interval.

\subsection{Persistence with the same cap core}

Fix $x$ with $R_T(x)>r_0^{-2}$. At sufficiently late regular times it
remains above threshold. Select recurrent neighborhoods in one of the
four canonical alternatives. Scalar and diameter controls, followed by
the temporal estimates, capture them in one compact regular set. All
metric jets converge there. Compact positive and round components retain
their strict controls after enlargement. For a strong neck retain a
smaller spatial carrier and backward interval; actual worldlines agree on
overlaps by uniqueness. Limiting normalized jets give the terminal neck.

A cap requires more than convergence of its underlying set. The selected
point must remain in the same core; its boundary must still be a central
sphere; all scalar, size and analytic estimates must hold on the resulting
carrier. Choose a sufficiently late recurrent cap in the common compact
regular set. Put $L=\varepsilon^{-1}$, $\alpha=c_T\varepsilon$, and
$r=\alpha^{-1}\geq200$. In its old end coordinate $s\in(-L,L)$,
translate the new terminal end chart by $-L+r$ and use $(-r,r)$.
Its image is the shortened end
\[
                         -L<s<-L+2r.
\]
Keep the old core literally fixed. Shorten the boundary-centered neck
without moving its central sphere. The translated end meets that same
sphere and core. Its negative quarter is $-L<s<-L+r/2$, retaining the
collar required between the compact core and end.

Compact metric convergence compares terminal and old normalizations. The
quantitative cap input is applied on the set

\[
 U=\text{(closed original core)}\;\cup\;
   \text{(the truncated end region)}.
\]

The whole set $U$ lies in the common compact regular set used for convergence.
The source estimates give on $U$ a positive
scalar lower bound, a scalar ratio bound, a gradient bound, an evolution bound,
and an upper bound for its total volume. Core balls retain their separate
volume lower bounds. The intrinsic-diameter estimate is applied to the
same truncation, while the core-ball calibration supplies the strict radius
$3/(4C)$; this is the margin that survives after the terminal cap constant is
doubled. The boundary neck is recentered first, and the translated end neck
is then restricted to the interval above. Consequently the old core point,
the boundary central sphere, and the complete truncated carrier all belong to
one terminal cap certificate. No estimate is inferred from convergence of the
underlying core alone.

The full-neck scalar estimate \eqref{eq:neck-cap-full-scalar} has coefficient
$10^{12}$. If $Q$ is the terminal scalar at the old center and $R$ that
at the translated center, write $c=R/Q$. The estimates give
$0<c\leq2$ and $|c-1|\leq S\varepsilon$, where
$S=2000000000004$. On the full old coordinate domain the squared jet
error is at most $4\varepsilon^2$ through order
$\lfloor(2\varepsilon)^{-1}\rfloor$. Translation has no derivative
cost, and changing normalization by $c$ bounds the new squared error by
\[
 2c^2(4\varepsilon^2)+6(c-1)^2
 \leq(32+6S^2)\varepsilon^2.
\]
The target order is no larger because $2\varepsilon\leq c_T\varepsilon$.
Consequently the translated normalized jets satisfy the squared budget
\begin{equation}
                         32+6S^2<c_T^2.
 \label{eq:singular-cap-budget}
\end{equation}
This makes the terminal chart a neck of accuracy $c_T\varepsilon$.
The fixed core retains strict metric and scalar margins: its core-ball lower
coefficient is $3/(4C)>1/(2C)$, supplying the lower inclusion for $2C$ after
changing to the exact terminal scalar radius. The shortened end gives the
outer clearance and diameter bound. The captured volume and the gradient and
evolution estimates transfer on the whole truncated carrier, so the terminal
certificate has the required scalar, jet, diameter, and volume fields with
the same enlargement to $2C$. These are checks on the retained core and its
boundary, not just on a nearby neck.
The explicit budget explains the departure from the printed factor two
in \cite[Lemma 11.23 and Claim 11.24, pp. 281--282]{MT2007}.

\subsection{Ends of the terminal region}

Proper scalar bounded below tends to infinity along an escaping end.
Beyond a compact sublevel all points have terminal canonical neighborhoods.
The neck--cap encounter theorem orders their overlaps into tubes and cores.
A compact canonical component cannot occur in an escaping tail. A cap
core may close one direction, but cannot recur indefinitely in the chosen
escaping direction: the encounter alternatives would close that direction
or form a compact component. A remaining tail is a neck tube.

Choose one central sphere far along it. Directed overlap parameterizations
give a smooth outgoing product; collar straightening fixes this sphere
and makes it the boundary. The product is proper because a compact set
cannot contain points escaping to infinite scalar. Reparameterize its
outgoing coordinate to $[0,1)$. This is the strong horn of
Definition~\ref{def:singular-terminal-horn}, with the actual terminal
strong necks already constructed. No completeness of $g(T)$ follows or
is needed \cite[Definitions 11.25--11.27 and Lemmas 11.28--11.30,
pp. 283--286]{MT2007}.

\section{Deep horns and strong-neck improvement}
\label{sec:singular-deep-horn}

A horn boundary is below $b>0$ when $R_T\leq b^{-2}$ there.
Write $\Omega_\rho=\{R_T\leq\rho^{-2}\}$ for the compact scalar sublevel.

\begin{theorem}[Uniform deep-horn height]
\label{thm:singular-m32}
\lean{PoincareMT.m32HornSelection}
\uses{thm:singular-m31,thm:singular-m30,prop:singular-m29}
There exists $0<\varepsilon_2\leq1/200$ allowing the following choices
in order. Fix $0<\varepsilon\leq\varepsilon_2$.
Choose positive parameters $r_0,C,C_a,\rho,\delta$, with $\rho<r_0$,
and a neck--cap topology theorem whose threshold contains
$c_T\varepsilon$. Then choose
\begin{equation}
                   0<h\leq\min\{\rho\delta,\rho/(2C)\}.
 \label{eq:deep-height}
\end{equation}
For every singular-time datum with these constants, every supplied
terminal conclusion, and every strong $\varepsilon_T=c_T\varepsilon$
horn $H$ whose boundary is below $\rho/(2C)$, each $x\in H$ with
$R_T(x)\geq h^{-2}$ centers a strong $\delta$-neck with terminal carrier
contained in $H$. There is such a neck with $R_T(x)=h^{-2}$ whose central
sphere cuts off a connected component of $H$ minus that sphere containing
a full parameter tail, escaping every compact set, and disjoint from
$\Omega_\rho$.
\source{MT2007}\dcref{Theorem 11.31 and Claim 11.35, pp. 287--292}
\end{theorem}

The supplied terminal conclusion is universally quantified. This theorem
does not choose another regular limit in place of the given extension.

\subsection{The first limit and its product direction}

If no height yields contained necks, select failures with terminal scalar
$Q_k\to\infty$ for fixed numerical data and rescale at the failing
points. The boundary scalar bound and boundary necks force every fixed
normalized base ball away from the horn boundary. Such balls also avoid
the missing end: otherwise unbounded scalar at bounded normalized distance
would contradict Theorem~\ref{thm:singular-bounded-distance}. Properness
of the horn and terminal scalar give compact closure of each ball.
Confined horn minimizers prevent shortcuts through the complement when
these actual terminal distances are compared with product coordinates. The
inverse convergence maps provide the needed quantitative confinement: for
every fixed source radius $r>0$ and $C>1$, the source ball of radius $r$
is eventually covered by the image of the limit ball of radius $Cr$.
The actual inverse sends that source ball into the enlarged limit ball.
An inverse derivative bound and compact closure of a larger limit ball
prevent a short source path from leaving this image. Thus the captured
separator and both attached tails have compact
closure in the source, and a minimizing segment cannot leave the horn while
being transferred to the limit.

Terminal strong necks give uniform noncollapse and backward survival.
Terminal analytic estimates follow by limits in reference flow boxes;
earlier exceptional times are handled by the dense-time condition.
The common controls hold, so Theorem~\ref{thm:singular-m30} gives a short
closed backward limit.

Both the distance to the boundary and to the missing end diverge after
normalization. Confined minimizing segments between these two sides
produce a minimizing line in the limit. Nonnegative curvature splits the
terminal metric along its Busemann coordinate. Central neck spheres are
essential in the horn: a sphere separating its two parameter ends cannot
bound a compact region on its outgoing side. Retaining this property
under compact convergence gives a fixed compact set $K$ separating the
two tails of the limiting line. Here is the separation argument. The source
central spheres lie in a fixed ball of radius $D$. The minimizing segments
between the two distant ends cross those spheres; their crossing points
provide bounded anchors for the limiting line $\gamma$. Put $B_*=2D+1$,
so that $d(p,\gamma(0))\leq B_*$, and let
$K=\overline B(p,B_*)$, which is compact by completeness. For
$T>2B_*+1$, both $\gamma(-T)$ and $\gamma(T)$ lie outside $K$.
If a path outside $K$ joined them, choose small path-connected endpoint
neighborhoods with compact closures outside $K$. Capture the path and these
closures in one convergence chart. The nearby endpoints of the truncated
source minimizing segments can be joined to the transferred path inside
these neighborhoods. Inverse confinement puts every preimage of a source
central sphere inside $K$. The resulting source path therefore avoids that
sphere while joining its opposite sides, contradicting separation in the
horn. Thus $K$ separates the line tails.

To prove the transverse factor compact, suppose it were noncompact.
Choose $q\in\Sigma$ outside the projection
of $K$, and choose heights $\pm L$ beyond the height range of $K$.
Path connectedness of $\Sigma$ joins each line endpoint horizontally to
$(q,\pm L)$ without meeting $K$; the vertical segment through $q$ also
avoids $K$. Concatenation contradicts the separating property. Thus there
is a compact connected surface $\Sigma$ and an actual product
$\Sigma\times\mathbb R$ with parallel unit gradient in the line
direction. Nullity evolution preserves this direction throughout
each known finite backward slab. One product map thus describes a surface
Ricci flow and a fixed Euclidean line. Nonflat normalization and the
scalar strong maximum principle give positive surface curvature at
interior times.

\subsection{Cap exclusion on a finite product}

Fix a compact portion of an interior slab. Its positive scalar lower
bound becomes arbitrarily large in the unnormalized source. At regular
times its points have canonical neighborhoods. Compact alternatives are
excluded by capturing a product segment longer than their diameter allows.
If a cap remains, enlarge the capture set to contain its entire closed
core and frontier and pull both back through the actual inverse convergence
map. Curvature-plane alignment makes projection of its neck boundary to
$\Sigma$ a local diffeomorphism. This requires comparison of Ricci on
all unit directions, not just comparison of scalar curvature: the axial
direction in the product is Ricci-null, while the tangential directions
of a small neck have a uniform positive Ricci lower bound. A tangent
vector annihilated by projection would violate these two bounds.

The projection is one-sheeted by order in the real coordinate. Its image
is open and compact, hence all of the connected surface $\Sigma$.
In each finite fiber select the lowest height. These selected points
form a nonempty closed subset of the sphere. They also form an open
subset: in an inverse neighborhood of a selected point the other compact
fiber branches have a strict height gap, preserved after shrinking the
neighborhood. Connectedness makes every point lowest in its fiber.
Injectivity of the full position-height map then forces each fiber to
have one point. Thus the actual boundary is a graph $s=f(z)$.

No compact subset with nonempty interior in $\Sigma\times\mathbb R$
can have its entire frontier in that graph. Choose an interior point off
the graph and follow its vertical half-line away from the graph. To leave
the subset it must cross the frontier, which that half-line never meets.
It remains inside for an unbounded interval, contradicting compactness.
This works for both allowed cap cores; it does not assume every cap core
is a ball. Thus the relevant regular neighborhoods are strong necks.
The full core must be captured before applying the inverse map.

\subsection{The signed estimate and backward continuation}

The metric jets of a sufficiently accurate normalized neck bound
the scalar Laplacian by $|\Delta R|\leq R^2/6$. In dimension three
$|\operatorname{Ric}|^2\geq R^2/3$, whence
\begin{equation}
 \partial_tR=\Delta R+2|\operatorname{Ric}|^2\geq\tfrac12R^2>0.
 \label{eq:singular-positive-scalar-time}
\end{equation}
This is the sign required in Claim 11.35; the absolute derivative bound
does not imply it. On a compact time window apply it on regular pieces
and use continuity at the finitely many exceptional times. Scalar on
retained worldlines is bounded backward by its terminal value.

Let $B$ bound the first limit's terminal scalar. Choose a positive step
$c$ small enough in terms of $B$ and the analytic constants, before any
compact capture set or current horizon. At a regular time just inside
a finite stage's left edge, the preceding argument gives actual centered
strong necks on the capture set. The source construction is tested on a
relatively compact open set $U$ whose closure lies in a larger open set $W$; the
attachment neck is supplied at every point of $\overline U$. Its native
backward time is at least a half-normalized unit because
$c\leq(4M)^{-1}$ and the scalar bound is $R\leq Mq$. Hence the attached
backward cylinders extend beyond the edge. They agree with the existing
cylinder on overlap: their points agree at attachment and uniqueness
propagates agreement along vertical trajectories. Scalar comparison and the
guarded ODE give one scalar bound on the glued cylinder; pinching gives a
common full curvature bound and vanishing negative defect.

The extension retains noncollapse. Strong-neck volume estimates apply on
controlled parabolic balls and agree with the old estimates on overlaps; the
same-clock identity identifies a source trajectory with the glued map,
independently of which local neck patch supplied it. Thus pointwise
scale-bounded noncollapse holds throughout the new cylinder, not merely at
time zero. The step is made on the original sequence:
every subsequence has a geometric stage limit on which the compact-product
argument constructs attachments. If good extended cylinders were not
eventually available originally, a subsequence of failures would contradict
this. Induction gives eventual noncollapsed cylinders of durations $nc$,
with one scalar bound chosen before $n$ and every spatial radius.
For any finite $T$, choose $n$ with $T<nc$ and restrict to $(-T,0]$.
These are the infinite-horizon long controls, so
Theorem~\ref{thm:singular-m30} now supplies an ancient limit.

The same scalar bound passes to this limit: its convergence maps and stage
cylinders agree at zero, hence along overlapping worldlines. Scalar
convergence bounds every point at every past time without identifying
unrelated stage limits. Splitting along the retained line gives an ancient
surface solution with positive noncollapse constant (the construction
retains $\kappa/2$). Theorem~\ref{thm:surface-ancient} makes
it round. A product slab with transverse factor $\mathbb{RP}^2$ is
nonorientable, whereas its retained embedding would make it an open subset
of the orientable horn $S^2\times[0,1)$. This excludes that alternative
and leaves the spherical factor. The limit is the shrinking round cylinder.

Choose a compact cylindrical segment long enough for a strong
$\delta$-neck, including its entire closed normalized backward interval.
Mixed-jet convergence, exact basepoint, and scalar normalization pull it
back to centered strong necks. The earlier ball capture keeps their terminal
carriers in the original horns, contradicting the chosen failures.

\subsection{Exact levels and an end cut}

First produce a seed cut without prescribing its exact scalar. Properness
captures the horn boundary and its intersection with $\Omega_\rho$ in
one compact connected parameter prefix $P$. Let $m$ be the maximum of
$R_T$ on $P$. Choose a center so far out that its scalar exceeds both
the requested level and $2m$. Neck scalar comparison puts its whole neck
above $m$, so it avoids $P$ and the horn boundary. Connectedness then
confines the neck to the horn. Its central sphere separates the outgoing
tail, and the horn nonfilling theorem excludes a compact outgoing side.
The resulting end cut avoids all of $P$, hence all of $\Omega_\rho$.

To move this cut to an exact level, put
\[
 b=\rho/(2C),\qquad q=2\max\{\rho^{-2},b^{-2}\}.
\]
In choosing $h$ impose the stronger
buffer $h\leq\min\{\rho/2,b/2\}$, so $q<h^{-2}$. The boundary and
forbidden sublevel are below $q/2$. Follow a fixed horn coordinate ray
from its boundary to a seed center above $h^{-2}$ and take its last
crossing of $h^{-2}$. The remainder is connected and stays above
$h^{-2}>q$; the contained-neck assertion applies at the crossing.

An end cut propagates through this connected high-scalar set. For nearby
centers, the sphere transport of Lemma~\ref{lem:neck-cap-sphere-isotopy}
and neck overlap supply a compactly supported diffeomorphism carrying
one central sphere to the other. Scalar comparison keeps its support above
$q/2$, disjoint from the forbidden sublevel and boundary, and confinement
keeps it inside the horn. It therefore carries the outgoing component,
full-tail and avoidance properties with it. The existence of such a cut
is locally constant in the center; connectedness propagates it from the
seed to the crossing, without any continuous choice of necks. This gives
the exact-level cut on the same terminal extension
\cite[Theorem 11.31, pp. 291--292]{MT2007}.

\section{A monotone scale on the same extension}
\label{sec:singular-selector}

\begin{corollary}[Uniform smaller heights]
\label{cor:singular-selector}
\lean{PoincareMT.M32DeepHornScaleSelection,PoincareMT.m32HornSelection}
\uses{thm:singular-m32}
For fixed sufficiently small $\varepsilon$, positive $C,C_a$, and the
supplied neck--cap topology theorem, there is a positive function
$h(\rho,\delta)$ on the positive quadrant, nondecreasing in each
variable, bounded above by $\min\{\rho\delta,\rho/(2C)\}$.
For every $r_0>\rho$ the full conclusion of
Theorem~\ref{thm:singular-m32}, including its exact-level neck and end
cut, holds at every $0<a\leq h(\rho,\delta)$ on every supplied terminal
extension.
\source{MT2007}\dcref{Corollary 11.36, p. 292}
\dcref{corrected index selection}
\end{corollary}

\begin{proof}
First obtain an interval of good heights. Put $b=\rho/(2C)$,
$r=\min\{\rho/2,b/2\}$, and choose $\eta\leq\delta$ below the
universal cut-propagation accuracy. Apply the deep-horn theorem at cutoff
$r$ and accuracy $\eta$, obtaining a seed cut at height $h_0$. Its
outgoing component has scalar greater than $r^{-2}$ and escapes to
infinite scalar. For $0<a<h_0$, continuity on that component and its
boundary center gives a point at exact scalar $a^{-2}$. The contained-neck
conclusion applies since $a^{-2}>h_0^{-2}$.

The choice of $r$ puts both the forbidden sublevel and horn boundary below
$r^{-2}/2$. On the connected set consisting of seed center and outgoing
component scalar stays above $r^{-2}$. Neck overlap and separation
propagate the outgoing cut along this set, preserving avoidance of the
lower sublevel. Apply this at the new center, then restrict accuracy from
$\eta$ to $\delta$. The full conclusion holds for every
$0<a\leq h_0$; monotonicity of $R\geq a^{-2}$ alone would not supply
the new exact-level cut.

Remove dependence on the original canonical radius by restricting the
singular datum from $r_0$ to $\rho$ and applying the preceding construction
with cutoff $\rho/2$. This strengthens the scalar guard and preserves the
actual flow, reference maps, terminal metric and horn. It yields a positive
interval uniform over all $r_0>\rho$. The full good-height property is
upward closed in $\rho,\delta$: a larger cutoff asks to avoid a smaller
scalar sublevel, its boundary hypothesis is stronger, and larger accuracy
is obtained by restricting the neck chart.

Set $d_n=2^{-n}$. Choose $b_n>0$ such that every $0<a\leq b_n$ is good
at $(d_n,d_n)$, shrink it to satisfy
$b_n\leq\min\{d_n^2,d_n/(2C)\}$, and put
$c_n=\min_{j\leq n}b_j$. For positive
$\rho,\delta$ define
\[
 n(\rho,\delta)=\min\{n:d_n\leq\rho\text{ and }d_n\leq\delta\},
 \qquad h(\rho,\delta)=c_{n(\rho,\delta)}.
\]
The eligible set is nonempty. Increasing either parameter decreases its
least index; since $c_n$ is nonincreasing, $h$ is nondecreasing. Positivity
comes from a finite minimum. Every $a\leq h$ is good at the selected
dyadic pair and hence at $(\rho,\delta)$, giving the upper bounds too.
The printed instruction to take a largest eligible index would fail,
since all sufficiently large indices are eligible
\cite[Corollary 11.36, p. 292]{MT2007}.
\end{proof}

The output is an extension of the supplied flow, its proper regular
terminal scalar, and contained strong necks with end cuts on that same
extension. Replacing the cut-off region by a cap, controlling its new
metric, and continuing the flow belong to the surgery chapter.
